\section*{Hoofdstuk \ref{ch:g11:the-eye} --- Het oog en zijn fotoreceptoren}

\begin{solution}{exo:g11:the-eye:1}
Hoornvlies, pupil (door de iris), ooglens, glasachtig lichaam, en daarna
de doorzichtige lagen van het netvlies (ganglioncellen en bipolaire
\omterm{def:g10:cells-common-unit:cell}{cellen}) tot aan de \omterm{def:g11:the-eye:photoreceptors}{fotoreceptoren} achterin.
\end{solution}

\begin{solution}{exo:g11:the-eye:2}
\omterm{def:g11:the-eye:photoreceptors}{Staafjes}: ongeveer 120 miljoen, heel gevoelig, één kleurstof, nachtzicht
in grijs. \omterm{def:g11:the-eye:photoreceptors}{Kegeltjes}: ongeveer 6 miljoen, hebben fel licht nodig, drie
kleurstoffen, zicht bij daglicht en kleurenzien, het scherpst in de gele
vlek.
\end{solution}

\begin{solution}{exo:g11:the-eye:3}
De gele vlek is het centrale kuiltje van het netvlies, dat alleen uit
dicht opeengepakte \omterm{def:g11:the-eye:photoreceptors}{kegeltjes} bestaat, waar het zicht het scherpst is. De
blinde vlek is de plaats waar de oogzenuw vertrekt, zonder
\omterm{def:g11:the-eye:photoreceptors}{fotoreceptoren}.
\end{solution}

\begin{solution}{exo:g11:the-eye:4}
\omterm{def:g11:the-eye:photoreceptors}{S-kegeltjes} ongeveer 45\%, \omterm{def:g11:the-eye:photoreceptors}{staafjes} ongeveer 60\%, \omterm{def:g11:the-eye:photoreceptors}{M-kegeltjes} ongeveer
20\%, \omterm{def:g11:the-eye:photoreceptors}{L-kegeltjes} ongeveer 10\%.
\end{solution}

\begin{solution}{exo:g11:the-eye:5}
Op het \omterm{def:g11:cell-cycle-mitosis:chromosome}{X-chromosoom}; de aandoening is \omterm{def:g11:genes-and-disease:genetic}{recessief} en op de X gelegen: zij
komt tot uiting bij elke man die het \omterm{def:g10:universal-dna:gene}{allel} draagt, en bij vrouwen alleen
met twee kopieën, en dus veel vaker bij mannen.
\end{solution}

\begin{solution}{exo:g11:the-eye:6}
'S Nachts reageren alleen de \omterm{def:g11:the-eye:photoreceptors}{staafjes}, en die hebben één enkele
kleurstof: zij melden helderheid, geen golflengte, dus geen kleur. De
gele vlek heeft alleen \omterm{def:g11:the-eye:photoreceptors}{kegeltjes}, die te ongevoelig zijn voor de ster;
kijk je opzij, dan valt het beeld van de ster op \omterm{def:g11:the-eye:photoreceptors}{staafjes}, die haar wel
opmerken.
\end{solution}

\begin{solution}{exo:g11:the-eye:7}
Dezelfde kleur: de verhoudingen S:M:L zijn in beide gevallen 0:0.6:1. Een
andere helderheid: licht 2 prikkelt elk \omterm{def:g11:the-eye:photoreceptors}{kegeltje} half zo sterk. Kleur is
de verhouding; helderheid is het totaal.
\end{solution}

\begin{solution}{exo:g11:the-eye:8}
Zonen krijgen hun X van hun moeder: allemaal normaal. Dochters krijgen de
X van hun vader: allemaal draagster, met normaal zicht. De zonen van elke
dragende dochter zijn met kans een half kleurenblind.
\end{solution}

\begin{solution}{exo:g11:the-eye:9}
Een man heeft één gemuteerde X nodig, een vrouw twee. Is de frequentie
van het \omterm{def:g10:universal-dna:gene}{allel} 0.08, dan draagt een vrouw er twee met kans
$0.08^2 = 0.0064$, ongeveer 0.6\% --- zestien keer minder dan 8\%.
\end{solution}

\begin{solution}{exo:g11:the-eye:10}
Verschillen hopen zich op met de tijd sinds een verdubbeling: 4\% tussen
L en M betekent een recente kopie; 57\% tussen elk van beide en S een veel
oudere. De verdubbeling L/M is de jongste gebeurtenis, de splitsing van S
ligt veel verder terug.
\end{solution}

\begin{solution}{exo:g11:the-eye:11}
De voorouder van de apen van de Oude Wereld verdubbelde zijn enige
opsinegen op de X; één kopie verschoof haar top naar langere golflengten
en gaf zo een derde type \omterm{def:g11:the-eye:photoreceptors}{kegeltje} en een nieuwe dimensie van kleur. Zij
bleef behouden, waarschijnlijk omdat zij hielp rijp fruit en jong blad
tegen het gebladerte te vinden.
\end{solution}

\begin{solution}{exo:g11:the-eye:12}
Blauw en geel (en blauwgroen van groen). Het \omterm{def:g10:universal-dna:gene}{S-gen} ligt op een gewoon
\omterm{def:g11:cell-cycle-mitosis:chromosome}{chromosoom}, dus moeten beide \omterm{def:g10:universal-dna:gene}{allelen} bij beide geslachten gemuteerd zijn
--- zeldzaam en onafhankelijk van het geslacht.
\end{solution}

\begin{solution}{exo:g11:the-eye:13}
$\num{2.5e-6}/\num{17e-3} \approx \num{1.5e-4}\,\unit{rad}$, ongeveer een
halve boogminuut. Op \qty{10}{m}: $\num{1.5e-4} \times 10 =
\qty{1.5}{mm}$.
\end{solution}

\begin{solution}{exo:g11:the-eye:14}
Een mannetje heeft één X, dus één \omterm{def:g10:universal-dna:gene}{allel}, dus één \omterm{def:g11:the-eye:photoreceptors}{kegeltje} van het
M/L-type: samen met S twee typen. Een vrouwtje met twee verschillende
\omterm{def:g10:universal-dna:gene}{allelen} brengt het ene in sommige \omterm{def:g11:the-eye:photoreceptors}{kegeltjes} tot uiting en het andere in
andere: drie typen, en dus zicht met drie kleuren; een vrouwtje met twee
gelijke \omterm{def:g10:universal-dna:gene}{allelen} heeft twee typen, net als een mannetje.
\end{solution}

\begin{solution}{exo:g11:the-eye:15}
De ganglioncellen en de oogzenuw zijn gaaf, dus bereiken elektrische
pulsen die aan hen worden gegeven de visuele hersenen als signalen. De
kwaliteit wordt begrensd door het aantal elektroden (honderden, tegen een
miljoen vezels en miljoenen zintuigcellen) en door het verlies van het
eigen sorteren van het beeld door het netvlies.
\end{solution}

\begin{solution}{pb:g11:the-eye:1}
\textbf{1.} \qty{470}{nm}: S 40, M 35, L 18. \qty{530}{nm}: S 0, M 100,
L 80. \qty{590}{nm}: S 0, M 40, L 90. \qty{650}{nm}: S 0, M 2, L 12.

\textbf{2.} \qty{470}{nm} prikkelt alle drie; \qty{650}{nm} in wezen
alleen L.

\textbf{3.} Met alleen S en M: \qty{590}{nm} geeft (0, 40) en
\qty{650}{nm} geeft (0, 2); \qty{530}{nm} geeft (0, 100). De lichtbronnen
van \qty{590}{nm} (oranje) en \qty{530}{nm} (groen) verschillen alleen in
de sterkte van M, als een zwak en een fel groen, en vormen het lastige
paar; \qty{650}{nm} lijkt op een heel zwak licht.

\textbf{4.} Met alleen S en L: \qty{530}{nm} geeft (0, 80) en
\qty{590}{nm} geeft (0, 90) --- vrijwel gelijk: groen en oranje worden
verward.

\textbf{5.} Beiden merken rood licht nog steeds op, met het type \omterm{def:g11:the-eye:photoreceptors}{kegeltje}
dat zij behouden; wat verloren gaat is de vergelijking tussen de reacties
van M en L, en juist die scheidt rood van groen.

\textbf{6.} $X^NX^c$: haar vader gaf haar zijn enige X, die het gemuteerde
\omterm{def:g10:universal-dna:gene}{allel} draagt; haar normale zicht laat zien dat de andere X normaal is.

\textbf{7.} Zonen: $X^NY$ normaal of $X^cY$ kleurenblind, elk de helft.
Dochters: $X^NX^N$ of $X^NX^c$, allebei met normaal zicht, elk de helft
(de tweede draagster).

\textbf{8.} Ines is draagster met kans $1/2$. Een dochter is kleurenblind
als zij van beide ouders $X^c$ krijgt: van de vader zeker, van Ines met
kans $1/2 \times 1/2$: $1/4$.

\textbf{9.} Een kleurenblinde dochter heeft twee gemuteerde
\omterm{def:g10:universal-dna:information}{X-chromosomen}, één van elke ouder; een vader geeft zijn enige X aan elke
dochter door, dus is zijn X gemuteerd en is hij kleurenblind.

\textbf{10.} Ja: haar zoon kreeg zijn X van haar, en die draagt het
gemuteerde \omterm{def:g10:universal-dna:gene}{allel}; met normaal zicht was zij $X^NX^c$.

\textbf{11.} Rodopsine takt het eerst af; daarna scheidt S zich van de
lijn L/M af; L en M splitsen als laatste: (rodopsine, (S, (L, M))).

\textbf{12.} 4\% verschil voor 35 miljoen jaar; 57\% zou ongeveer
$35 \times 57/4 \approx 500$ miljoen jaar kosten --- in overeenstemming
met de vroege oorsprong bij de \omterm{def:g10:common-ancestry:bodyplan}{gewervelden} (het verband is bij grote
verschillen niet precies evenredig, dus dit is een orde van grootte).

\textbf{13.} De twee vrijwel gelijke, naast elkaar liggende \omterm{def:g10:universal-dna:gene}{genen} kunnen
tijdens het paren van de twee \omterm{def:g10:universal-dna:information}{X-chromosomen} verkeerd uitlijnen; een
uitwisseling levert dan de ene X met drie kopieën en de andere met één
kopie op. Een man die de X met één kopie krijgt, heeft maar één \omterm{def:g11:the-eye:photoreceptors}{kegeltje}
van het M/L-type: rood-groenkleurenblindheid.

\textbf{14.} Het paren berust op gelijkenis in volgorde: twee naast elkaar
liggende \omterm{def:g10:universal-dna:gene}{genen} die voor 96\% gelijk zijn, kunnen uit de pas paren, \omterm{def:g10:universal-dna:gene}{gen} 1
met \omterm{def:g10:universal-dna:gene}{gen} 2, en ongelijk uitwisselen.

\textbf{15.} Zij trad één keer op, bij de gemeenschappelijke voorouder van
de apen van de Oude Wereld, na hun scheiding van de andere zoogdieren en
van de apen van de Nieuwe Wereld, zo'n 35 miljoen jaar geleden, en werd
door al hun nakomelingen geërfd.

\textbf{16.} $\num{2.5e-6}/\num{17e-3} \approx \num{1.5e-4}$ rad, ongeveer
$0.5'$.

\textbf{17.} $\num{1.5e-4} \times 6 \approx \qty{0.9}{mm}$: fijner dan de
streepjes van \qty{1.7}{mm}, en daarom leest normaal zicht die regel
moeiteloos (in de praktijk brengt optische onscherpte de grens dicht bij
$1'$).

\textbf{18.} Minstens tien keer grover (100 \omterm{def:g11:the-eye:photoreceptors}{staafjes} delen één signaal,
ongeveer 10 bij 10): enkele boogminuten; in ruil daarvoor bundelen honderd
zintuigcellen hun licht en merken zij veel zwakkere bronnen op.

\textbf{19.} Twee keer de dichtheid betekent een $\sqrt{2}$ keer kleinere
afstand, en een groter \omterm{def:g11:the-eye:parts}{oog} voegt een langere brandpuntsafstand toe:
ongeveer twee keer zo fijn.

\textbf{20.} De S-, M- en \omterm{def:g11:the-eye:photoreceptors}{L-kegeltjes}, die elk vrijwel blind zijn voor het
uiteinde van het spectrum dat het verst van hun top ligt; één verdubbeling
naast elkaar van het opsinegen op de X, 35 miljoen jaar geleden, gaf de
primaten het derde; de gele vlek onderscheidt ongeveer een halve
boogminuut.
\end{solution}
